\documentclass[10pt]{article}

\usepackage[utf8]{inputenc}

\usepackage[T1]{fontenc}

\usepackage{amsmath}

\usepackage{amsfonts}

\usepackage{amssymb}

\usepackage[version=4]{mhchem}

\usepackage{stmaryrd}

\usepackage{bbold}



\begin{document}

Системы уравнений (1) и (2) позволяют, зная закон движения тела, определить момент действующих на него сил и наоборот, зная действующие на тело силы, огределить закон его движения.



Лuт.: [1] Бу хгольц Н. Н., Основной курс теоретической механики, 6 изд., ч. 2, М., $1972 . \quad$ C. М. Таре. ЭЙЛЕРА ФОРМУЛА - см. Динамика твердого тела.



эйЛЕРА - КРЫЛОВА УГЛЫ - то же, что ко р а бе л ь ны е углы (см. Ориентации уалы).



ЭЙЛЕРА - ЛАГРАНХА УРАВНЕНИЯ - см. Вариационное исчисление, Вариационное исчисление в целом, Лагранжев формализм.



ЗЙЛЕРА - ПУАНСО СЛУЧАЙ - первыЙ (по времени и простейший) из трех классических случаев интегрируемости уравнений движения тяжелого твердого тела около неподвижной точки (другие - Лагранжа - Пуассона случай, Ковалевской случай). Он характеризуется тем, что неподвижной точкой является центр тяжести тела (обычно совпадающий с началом координат $O$ ). В этом случае $x_{0}=y_{0}=z_{0}=0$ (см. Твердое тело; движение около неподвижной точки) и уравнения Эйлера - Пуассона (I)-(II) обладают следующим четвертым дополнительным первым интегралом:



\[

A^{2} p^{2}+B^{2} q^{2}+C^{2} r^{2}=l^{2}=\mathrm{const}

\]



выражающим сохранение кинетич. момента; открыт Л. Эйлером [1].



Этот интеграл и позволяет свести интегрирование уравнений (I)-(II) в рассматриваемом случае к квадратурам, обращению нек-рого эллиптич. интеграла (2-го рода) и решению квадратных уравнений. В частном случае, когда начальные данные таковы, что $B=\frac{l^{2}}{h}$, функции $p, q, r$ выражаются через элементарные функции (см. [3]). Решение уравнений II суть



\[

\gamma=\frac{A p}{l}, \quad \gamma^{\prime}=\frac{B q}{l}, \quad \gamma^{\prime \prime}=\frac{C r}{l} .

\]



Если $p, q, r$ известны как функции $t$, то легко определить и положение твердого тела, применяя выражение $\gamma, \gamma^{\prime}, \gamma^{\prime \prime}$ через Эйлера углы.



Заслуга Л. Пуансо (см. [2]) состоит в том, что он дал геометрич. интерпретацию движения твердого тела в рассматриваемом случае.



Первое представление Пуансо: эллипсоид инерции (уравнение к-рого $A x^{2}+B y^{2}+C z^{2}=1$ в репере, связанном с телом) катится без проскальзывания по неподвижной плоскости $\pi$, перпендикулярной кинетич. моменту $l$ и отстоящей от неподвижной точки $O$ на расстояние



\[

\delta=\sqrt{2 h} /|l|

\]



Эллипсоид, катаясь по плоскости $\pi$, оставляет на ней след, называемый ге р п о л о д и е й; точка касания с плоскостью на эллипсоиде описывает кривую, называемую поло д и е й. Так как полодия высекается на эллипсоиде инерции направлением угловой скорости $\omega$, то полодии совпадают с фазовыми кривыми уравнений Эйлера при



\[

T=\frac{1}{2}\left(A p^{2}+B q^{2}+C r^{2}\right)

\]



Герполодии располагаются на плоскости в нек-ром кольце и, вообще говоря, не замкнуты.



Второе представление Пуансо: если с подвижным репером (связанным с телом) связать конус



\[

\frac{A \xi^{2}}{A-D}+\frac{B \xi^{2}}{B-D}+\frac{C \xi^{2}}{C-D}=0, \quad D=\frac{l}{h},

\]



то движение тела можно представить как качение этого конуса по плоскости $\pi$, проходящей через $O$ перпендикулярно кинетич. моменту $l$ и врашаюшейся с постоянной угловой скоростью.



Если тело обладает динамич. симметрией, то эллипсоид Пуансо становится эллипсоидом вращения и движение можно описать качением прямого кругового конуса, фик- сированного в теле, по неподвижному в пространстве прямому круговому конусу.



Jum.: [1] Eule r L., "Mémoires de l'Acad. des Sc. de Berlin», 1758 (см. также пер.: Э йл ер Л., Основы динамики точки, М., 1938); [2] Poinsot L., «J. de Liouville", 1851, t. 16; [3] Голубе в В. В., Лекции по интегрированию уравнений движения тяжелого твердого тела около неподвижной точки, М., 1953; [4] Т а т а р и но в Я. В., Лекции по классической динамике, М., 1984. М. И. Войцеховский. ЭЙЛЕРА - ПУАССОНА УРАВНЕНИЯ - сМ. Твердое тело; движение около неподвижной точки.



ЭЙЛЕРОВА ХАРАКТЕРИСТИКА - топологический инвариант фигуры Ф, определяемый с помощью разбиения ее на клетки разных размерностей; кл е т ка - гомеоморфный образ куба, под ф и гу р о й понимается конечное объединение всевозможных клеток или (более общо) гомеоморфных им выпуклых многогранников, быть может, вырожденных. Тогда эйл е р о в а х а р а к т р и с т и ка - это число $\chi(\Phi)$, поставленное в соответствие каждой фигуре так, что выполнены аксиомы:



\begin{enumerate}

  \item Э. Х. пустой фигуры равна нулю: $\chi(\varnothing)=0$;



  \item для каждого (в том числе вырожденного) выпуклого многогранника $C$ Э. х. равна единице: $\chi(C)=1$;



  \item для любых фигур $A$ и $B$ имеет место аддитивность:



\end{enumerate}



\[

\chi(A \cup B)+\chi(A \cap B)=\chi(A)+\chi(B) \text {. }

\]



Так, для многогранника в трехмерном пространстве Э.х. выражается формулой:



\[

\chi=\mathrm{B}-\mathrm{P}+\Gamma

\]



где В - число вершин, Р - число ребер, Г - число граней; для замкнутой поверхности в $\mathbb{R}^{3}$ Э. х. выражается через гауссову кривизну формулой Гаусса - Бонне:



\[

\int K d s=2 \pi \xi

\]



М. И. Войцеховский.



ЭЙНШТЕЙНА ПРИНЦИП ОТНОСИТЕЛЬНОСТИ, прин цип относ ительности: все физические явления (механические, электромагнитные и любые другие) при одинаковых начальных условиях протекают одинаково во всех инерииальных системах отсчета. М. И. Войцеховский. ЭЙНШТЕЙНА ТЕНЗОР - тензор вида



\[

G^{i j}=R^{i j}-\frac{1}{2} g^{i j} R, \quad R=g^{i j} R_{i j},

\]



где $R^{i j}=g^{i l} g^{j m} R_{l m}$ и $R_{i j}-Р$ так наз. уравнениям Эйнштейна



\[

G^{i j}, j=0

\]



в пустоте, описывающим гравитационное поле (метрику) в пустом пространстве (без материи). ЭЙНШТЕЙНА УРАВНЕНИЯ - основные уравнения общей теории относительности, связывающие геометрическую структуру пространства-времени и тензор энергии-импульса всевозможных полей и частиц, находящихся в нем. Э. у. имеют вид:



\[

R_{i k}-\frac{1}{2} g_{i k} R=\frac{8 \pi G}{c^{4}} T_{i k}, \quad i, k=0,1,2,3,

\]



где $R_{i k}-$ тензор Риччи, $R$ - скалярная кривизна, $T_{i k}-$ тензор энергии-импульса, $G$ - гравитационная постоянная, $c$ - скорость света, $g_{i k}-$ метрич. тензор пространства-времени. В Э.у., в отличие от других уравнений математич. физики, ищутся одновременно и неизвестные компоненты метрич. тензора и то многообразие. на к-ром они заданы. Э. у. содержат 1-е и 2-е производные от компонент метрич. тензора, следствием чего являются уравнения движения среды $T_{i k}^{k}=0$. Наиболее разработана постановка задачи Коши для Э.у., в к-рой начальные данные задаются на нек-рой пространственноподобной гиперповерхности и ищется эволюция во времени этой гиперповерхности



ЭЙНШТЕЙНА 673



43 Математическая физика





\end{document}